\begin{abstract}
Which components of a source's earlier history are distinguishable from black-hole images when the spacetime is known but the rest of the source is not? We formulate this question as retarded-time tomography and separate historical reach, historical dimension, and held-out recovery. The central finite-noise result concerns a 72-dimensional old, non-axisymmetric target within a localized 224-dimensional source model: the sampled direct operator is exactly null on the target, whereas the order-labelled stack retains two operational combinations after all other 152 coefficients are profiled out. This count uses the frozen legacy measure; fixed-row reweighting bounds permit one or two. A support theorem explains the geometric mechanism: temporal factors disjoint from the direct retarded-time footprint generate zero direct columns, while added orders can produce nonzero responses. On a twelve-geometry grid at matched direct-reference $\SNR=100$, the first two indirect orders extend the anchor-connected detectable span by $40M$, $24M$, and one $4M$ grid step at inclinations $20^\circ$, $50^\circ$, and $75^\circ$. A separate three-anchor enrichment study finds increasing absolute supported count alongside decreasing supported fraction. Two sealed benchmarks then test different historical objects: baseline-inclusive emissivity span increases from $48M$ to $80M$ in the successful regimes, while median paired morphology-error reductions are $0.133$ for truncated SVD and $0.164$ for ridge. The level result is baseline dominated; severe mismatch and reliable multi-feature recovery remain negative, and the stable morphology interval is zero. Order attribution shows the two operational combinations are not driven by the numerically weakest image: omitting $n=2$ leaves both above threshold, with $96.90\%$ and $95.64\%$ of their residualized response energy in $n=1$. The results establish a finite-model historical-information and recovery analysis under ideal order labelling. Full transfer quadrature remains unqualified, and the index-sum control supplies no physical order-resolution attribution. The Mahakal phenomenon names the resulting separation between access to earlier source times and recoverable historical content.
\end{abstract}

\section{Introduction}\label{sec:intro}
Light from the vicinity of a black hole can reach a distant observer along substantially different null paths. Direct rays and increasingly bent higher-order rays sample different source positions and emission times. The resulting image hierarchy has established demagnification, rotation, and delay structure \cite{gralla,johnson}. For a time-dependent source, an image order is not a single delayed snapshot: its rays carry a distribution of delays, and different orders have overlapping retarded-time footprints. The inverse question is then precise: \emph{given the spacetime and an observation model, which parts of the earlier source are distinguishable even when the remaining emission is unknown?}

We answer this question at three levels. \emph{Historical reach} asks which earlier source times are sampled and where a prescribed localized perturbation exceeds a sensitivity threshold. \emph{Historical dimension} counts distinguishable directions in a declared source space at finite noise, including when other source coefficients are unknown. \emph{Historical recovery} tests a fixed estimator on unseen histories under a declared loss. Separating these levels is necessary because a sampled old perturbation can be weak, can be mimicked by another source component, or can be detectable without supporting a reliable sequence of reconstructed images. The hierarchy turns an appealing delayed-image picture into distinct mathematical and experimental questions.

\begin{definition}[The Mahakal phenomenon]
We use the term for the separation between gravitational access to earlier source times and the finite-model dimension and fidelity of the history supported by the observations. Higher-order footprints can enter time regions missed by the direct footprint; resolving that history in a richer source representation can expose null or weak directions hidden by a coarser description. The relation is geometry-, noise-, and source-class-dependent. It is not a new propagation law, a statement that source enrichment always destroys injectivity, or a claim that the black hole contains a recoverable movie.
\end{definition}

The main information result concerns a 72-dimensional old, non-axisymmetric target within a 224-dimensional localized source model. At the reference geometry, every direct target column is exactly zero. After allowing the other 152 coefficients to vary as nuisance parameters, the labelled stack retains two source-normalized combinations above the declared noise threshold. These nuisance coefficients include recent emission, old axisymmetric factors, and boundary-overlapping factors. The result therefore identifies historical responses that cannot be reproduced by adjusting the remainder \emph{within that model}; it does not assume that the remainder is known. The surviving combinations are temporally broad rather than independently reconstructed frames.

Two held-out studies then ask what fixed estimators recover. The emissivity-level study gives a longer stable span under its baseline-inclusive field metric. The morphology study gives material reductions in a resolution-aware, all-state error. Their distinct negative endpoints---level dominance, severe-mismatch failure, unreliable multiple features, and a zero stable morphology interval---locate the boundary between additional information and recovery of increasingly detailed historical objects. This boundary is a measured part of the result, not a reason to equate all forms of recovery.

\subsection{Scope and contribution}
The setting is a known, stationary Kerr or Schwarzschild geometry, an equatorial source, and a monochromatic scalar-intensity image-domain model. We retain orders $n=0,1,2$ under the ray-map classification. Multi-geometry sensitivity is evaluated on twelve prescribed spin--inclination pairs; source-class enrichment uses three prescribed anchors; held-out inversion uses one reference geometry. These are separate experimental populations.

We combine a support-nullity theorem, source-normalized nuisance projection, source-class stress, and predeclared held-out losses to evaluate source-history identifiability. The theorem establishes exact direct blindness under support disjointness; the nuisance calculation tests whether added-order responses remain distinct when other coefficients are free; the inverse benchmarks test estimator performance rather than inferring it from rank. \Cref{tab:evidence-map} summarizes the resulting evidence without merging its source spaces or metrics. The compiled experiment records and amendments supply the numerical results \cite{archive}.

\begin{table}[htbp]
\centering\small
\caption{Three questions, with different endpoints and evidence. All numerical entries refer to their declared finite operators and noise conventions. An improvement in one row does not establish the stronger endpoint in another.}
\label{tab:evidence-map}
\begin{tabularx}{\linewidth}{@{}p{0.15\linewidth}YY@{}}
\toprule
Question&Result&Interpretation and boundary\\
\midrule
Reach&Added-order anchor-connected span increases by $40M$, $24M$, and $4M$ across the three tested inclinations.&Detectability of the prescribed scalar probe, not an estimator's reconstructed interval.\\
Dimension&Two of 72 old-target combinations pass after profiling 152 other coefficients; direct target response is zero.&A source-normalized, legacy-measure result; fixed-row reweighting bounds allow one or two.\\
Recovery&Level span $48M\to80M$ in successful regimes; morphology reductions $0.133$ and $0.164$.&Different source objects and losses. Severe mismatch and multi-feature recovery remain negative; stable morphology span is zero.\\
\bottomrule
\end{tabularx}
\end{table}

\Needspace{4\baselineskip}
The empirical scope is a finite-model benchmark, not a completed observational demonstration. Continuous transfer-integral accuracy and the physical interpretation of an unresolved spatial image require separate evidence. We state these boundaries with the relevant results while preserving the positive conclusions at the scope their operators, source norms, and estimators support.

\section{Relation to existing work}\label{sec:related}
The Kerr image hierarchy and its long-baseline interferometric signatures provide the established propagation setting \cite{gralla,johnson}. Glimmer and photon-ring autocorrelation studies analyze delayed image structure and correlations across image position and time \cite{wong,hadar}. In the latter, correlation profiles contain source-statistical information as well as lens-dependent structure. We therefore do not identify the existence of delayed source information as the novelty of the present work.

Two particularly close developments help locate the inverse question. Wong et al.\ propose correlating a total light curve with high-spatial-frequency interferometric signals to estimate an echo delay in simulated accretion-flow observations \cite{echoes2024}. Rojas-Paternina and C\'ardenas-Avenda\~no analyze image-order delay distributions and their interaction with source variability, comparing slow-, fast-, and brisk-light propagation prescriptions \cite{propagation2026}. These studies establish direct precedents for echo observables and distributed retarded-time image formation. Here the geometry is fixed and the target is different: the source-history combinations that remain distinguishable after the modelled remainder is profiled out, followed by held-out recovery of specified historical objects.

Existing tomography and imaging methods also provide reconstruction precedents. Hotspot-based spacetime tomography studies constrained dynamical sources \cite{tiede}; dynamic interferometric imaging reconstructs evolving images using an explicit temporal model \cite{bouman}; orbital polarimetric tomography combines a gravitational model and an emission representation to infer flare structure \cite{levis}. The present comparison does not assert that these methods lack source information or that they are alternative implementations of the same benchmark. Our analysis instead keeps the historical target, its source norm, the unknown remainder, and the reconstruction loss explicit, so that support, likelihood information, and estimator performance can be evaluated separately.

AART supplies the analytic Kerr ray-tracing framework underlying the archived maps \cite{aart}. Accuracy of geodesic expressions is distinct from accuracy of a downstream detector integral built from sampling, masking, interpolation, and quadrature. The numerical qualification in \cref{sec:validation} accordingly separates pointwise transfer/domain checks, discrete operator mechanics, and full acquisition accuracy.

The broader distinction between identifiability and stable inversion is established in light-ray transforms and discrete inverse problems \cite{lassas,hansen}. Regularization and prior assumptions can stabilize or select a reconstruction, but they do not constitute additional likelihood information \cite{stuart}. The specific contribution here is the application of that distinction to the image-order retarded-time operator: an explicit direct-null historical target, its nuisance-adjusted finite-noise dimension, and bounded recovery tests. We make no literature-wide priority claim for tomography, photon echoes, or the general mathematics of inverse problems.
